\subsection*{Moduł 3: DLP (Atende DLP)}
\addcontentsline{toc}{subsection}{Moduł 3: DLP (Atende DLP)}

\subsubsection*{3.1. Cel systemu}

System \textbf{Atende DLP} (\emph{Data Loss Prevention}) jest rozproszonym rozwiązaniem do wykrywania
i przeciwdziałania nieautoryzowanemu wyciekowi danych wrażliwych ze stacji roboczych w organizacji.
System monitoruje operacje wykonywane przez użytkowników na plikach oraz w schowku systemowym,
przechwytuje kopiowane pliki i treści, a następnie analizuje je pod kątem obecności danych
wrażliwych (np. numerów PESEL, numerów telefonów, treści zgodnych z regułami zdefiniowanymi przez
administratora). Wykryte naruszenia są rejestrowane jako alerty, prezentowane w panelu
administracyjnym oraz -- w przypadku alertów krytycznych -- zgłaszane e-mailem odpowiedzialnym
osobom.

System składa się z trzech niezależnie rozwijanych i wdrażanych modułów, które razem tworzą
kompletny łańcuch: \emph{zbieranie danych na stacji roboczej} $\rightarrow$ \emph{gromadzenie
i przechowywanie po stronie serwera} $\rightarrow$ \emph{analiza treści i wykrywanie naruszeń}.

\subsubsection*{3.2. Moduły systemu -- przegląd}

\subhead{DLP Agent (\ep{dlp-agent})}
Usługa systemu Windows instalowana na stacjach roboczych użytkowników. Napisana w języku C\#
(.NET Framework 4.8), uruchamiana jako usługa Windows o nazwie \ep{AtendeDLPMonitor} przy użyciu
frameworku \textbf{Topshelf}. Agent w czasie rzeczywistym monitoruje operacje plikowe (poprzez
mechanizm \emph{Event Tracing for Windows}, ETW), śledzi sesje logowania użytkowników oraz uruchamia
towarzyszący proces monitorujący zawartość schowka systemowego (\ep{DLPClipboardMonitor.exe}).
Zebrane logi zdarzeń oraz skopiowane pliki są okresowo wysyłane do serwera backendowego przez HTTP/REST.
Agent \textbf{nie wystawia} własnego API -- jest wyłącznie klientem API serwera backendowego.

\subhead{DLP Backend Server (\ep{dlp-backend-server})}
Centralny serwer aplikacyjny napisany w Javie 17 z wykorzystaniem Spring Boot 3.4. Odpowiada za:
przyjmowanie logów i plików od agentów, przechowywanie ich w bazie danych MySQL/MariaDB,
udostępnianie API panelu administracyjnego (logowanie, przegląd logów, alertów, użytkowników,
grup, reguł DLP i raportów), a także za komunikację z lokalnym silnikiem parsująco-skanującym
(\ep{dlp-parser}) poprzez dedykowany, wewnętrzny interfejs HTTP. Serwer może pracować w jednym
z dwóch trybów (profili Spring): \textbf{\ep{collector}} (serwer wdrażany u klienta, zbiera dane
z agentów) oraz \textbf{\ep{aggregator}} (centralny serwer zbierający alerty z wielu instalacji
typu \emph{collector}). Jest to jedyny moduł systemu wystawiający publiczne REST API.

\subhead{DLP Parser (\ep{dlp-parser})}
Silnik analityczny napisany w Pythonie, uruchamiany jako zestaw skryptów wsadowych (nie jest
serwerem HTTP). Skrypty są cyklicznie wywoływane przez harmonogram (\emph{Script Scheduler})
wbudowany w serwer backendowy. Parser pobiera z backendu (poprzez wewnętrzne API) pliki oczekujące
na przetworzenie, ekstrahuje z nich tekst (parsery dokumentów Word/Excel/PDF, OCR obrazów),
a następnie uruchamia na wyekstrahowanej treści silnik reguł (parsery zdefiniowane w plikach JSON
oraz podpięte pod nie wtyczki Pythona) w celu wykrycia danych wrażliwych. Wynik analizy (lista
alertów) jest odsyłany z powrotem do backendu, a w przypadku alertów krytycznych -- rozsyłany
e-mailem.

\subsubsection*{3.3. Architektura i przepływ danych}

Poniższy diagram przedstawia relacje pomiędzy modułami systemu oraz kierunki komunikacji.

\begin{figure}[h]
\centering
\resizebox{\textwidth}{!}{%
\begin{tikzpicture}[
    node distance=1.6cm and 2.2cm,
    box/.style={draw, rounded corners, minimum width=3.6cm, minimum height=1.3cm,
                align=center, fill=blue!6, font=\small},
    dbbox/.style={draw, cylinder, shape border rotate=90, minimum width=2.2cm,
                  minimum height=1.2cm, aspect=0.3, fill=gray!12, font=\small, align=center},
    arrow/.style={-{Latex[length=2.5mm]}, thick}
]
\node[box] (agent) {Stacja robocza\\\textbf{DLP Agent}\\(C\#, Windows Service)};
\node[box, right=of agent, xshift=1.2cm] (backend) {\textbf{DLP Backend Server}\\(Java / Spring Boot)\\profil \emph{collector}};
\node[box, right=of backend, xshift=1.2cm] (parser) {\textbf{DLP Parser}\\(Python, skrypty wsadowe)};
\node[dbbox, below=of backend] (db) {MySQL /\\ MariaDB};
\node[box, above=2.0cm of parser, fill=orange!8] (aggregator) {\textbf{DLP Backend Server}\\profil \emph{aggregator}\\(centralny)};
\node[box, below=1.4cm of agent, fill=green!8] (admin) {Panel\\administracyjny};

\draw[arrow] (agent.east) -- node[above,font=\scriptsize,align=center]{REST/HTTPS\\logi, pliki, konfiguracja} (backend.west);
\draw[arrow] (backend) -- (db);
\draw[arrow] (backend.east) -- node[above,font=\scriptsize,align=center]{API wewnętrzne\\(port 8081)} (parser.west);
\draw[arrow] (parser.west) to[out=200,in=-20] node[below,font=\scriptsize]{wynik: alerty} (backend.south east);
\draw[arrow] (backend.north) |- node[above,font=\scriptsize, pos=0.25]{alerty (harmonogram)} (aggregator.west);
\draw[arrow] (admin.east) -- node[below,font=\scriptsize]{JWT / REST} (backend.south west);
\end{tikzpicture}%
}
\caption{Architektura systemu Atende DLP i kierunki komunikacji między modułami.}
\end{figure}

\noindent Kluczowe cechy architektury:
\begin{itemize}
  \item Serwer backendowy udostępnia \textbf{dwa niezależne punkty wejścia HTTP} na osobnych
        portach: port główny (domyślnie \ep{8080}) dla agentów i panelu administracyjnego oraz
        port wewnętrzny (domyślnie \ep{8081}, prefiks \ep{/internal/}) wyłącznie dla lokalnego
        procesu parsera. Rozdzielenie portów wymuszane jest przez filtr
        \ep{DlpInternalEndpointsFilter}.
  \item Agent i parser \textbf{nie komunikują się ze sobą bezpośrednio} -- cała wymiana danych
        odbywa się za pośrednictwem serwera backendowego, który pełni rolę centralnego magazynu
        stanu (pliki, logi, statusy przetwarzania, alerty).
  \item W topologii wielooddziałowej serwery typu \emph{collector} (po jednym na lokalizację
        klienta) okresowo wysyłają zagregowane alerty do jednego, centralnego serwera typu
        \emph{aggregator}, korzystając z klienta Feign i tokenu autoryzacyjnego przypisanego do
        danego klienta.
\end{itemize}

% =========================================================================
\subsubsection*{3.4. API wystawiane przez serwery}
\label{sec:dlp-api}
% =========================================================================

Jedynym modułem systemu wystawiającym interfejs API jest \textbf{DLP Backend Server}. Agent oraz
parser są wyłącznie konsumentami tego API. API podzielone jest na cztery grupy funkcjonalne,
różniące się mechanizmem uwierzytelniania i portem, na którym są dostępne (zob. rozdz.~\ref{ch:bezpieczenstwo}):

\begin{itemize}
  \item \textbf{API uwierzytelniania i panelu administracyjnego} (\ep{/auth}, \ep{/admin}) --
        port główny, uwierzytelnianie JWT, aktywne zawsze.
  \item \textbf{API agenta} (\ep{/device}, \ep{/log}, \ep{/upload}) -- port główny, uwierzytelnianie
        nagłówkami \ep{WinKey}/\ep{BiosSN}, aktywne w profilu \emph{collector}.
  \item \textbf{API agregacji alertów} (\ep{/alert}) -- port główny, uwierzytelnianie tokenem
        klienta, aktywne w profilu \emph{aggregator}.
  \item \textbf{API wewnętrzne} (\ep{/internal/...}) -- osobny port wewnętrzny, bez
        uwierzytelniania (izolacja na poziomie sieci/portu), używane wyłącznie przez
        \ep{dlp-parser}, aktywne w profilu \emph{collector}.
\end{itemize}

Wszystkie odpowiedzi API (poza pobieraniem plików binarnych) mają wspólną kopertę
\ep{ApiResult} zawierającą co najmniej pola \ep{errorCode} oraz \ep{message}, rozszerzaną
w konkretnych odpowiedziach o właściwe dane (pojedynczy obiekt lub listę wraz z \ep{totalCount}
dla wyników stronicowanych).

\subhead{Auth API -- \texorpdfstring{\code{/auth}}{/auth}}

Interfejs logowania i zarządzania sesją administratora panelu. Dostępny zawsze, niezależnie od
aktywnego profilu.

\begingroup\footnotesize
\begin{longtable}{M{1.4cm} M{2.6cm} M{5.3cm} M{5.8cm}}
\toprule
\rowcolor{tablehead!15}
\textbf{Metoda} & \textbf{Ścieżka} & \textbf{Body} & \textbf{Opis} \\
\midrule
\endhead
\httpmethod{POST} & \ep{/auth/login} & \ep{LoginDTO\{login, password\}} &
  Logowanie administratora. Zwraca parę tokenów JWT (\ep{accessToken}, \ep{refreshToken}) wraz z czasem ważności. \\
\httpmethod{POST} & \ep{/auth/refreshToken} & \ep{RefreshTokenRequest\{refreshToken\}} &
  Odświeżenie tokenu dostępowego na podstawie ważnego tokenu odświeżającego. \\
\httpmethod{POST} & \ep{/auth/logout} & -- (JWT w nagłówku \ep{Authorization}) &
  Unieważnia bieżącą sesję administratora (dodaje token do listy odwołanych). \\
\httpmethod{POST} & \ep{/auth/logoutAll} & -- (JWT w nagłówku) &
  Unieważnia wszystkie aktywne sesje danego administratora. \\
\httpmethod{POST} & \ep{/auth/recoverPassword} & \ep{PasswordRecoveryRequest\{login, email\}} &
  Generuje jednorazowy token resetu hasła i wysyła e-mail z linkiem. Zawsze zwraca sukces (nie ujawnia, czy konto istnieje). \\
\httpmethod{POST} & \ep{/auth/resetPassword} & \ep{ResetPasswordRequest\{login, email, token, newPassword, newPasswordConfirmation\}} &
  Ustawia nowe hasło po weryfikacji tokenu resetu; wylogowuje wszystkie aktywne sesje administratora. \\
\bottomrule
\caption{Endpointy \code{AuthController}.}
\end{longtable}
\endgroup

\subhead{Admin API -- \texorpdfstring{\code{/admin}}{/admin}}

Interfejs panelu administracyjnego. Wszystkie endpointy wymagają roli \ep{ADMIN}, wymuszanej
adnotacją \ep{@PreAuthorize} z wyrażeniem \ep{hasAuthority('ADMIN')}.

\begingroup\footnotesize
\begin{longtable}{M{1.4cm} M{4.2cm} M{4.6cm} M{4.1cm}}
\toprule
\rowcolor{tablehead!15}
\textbf{Metoda} & \textbf{Ścieżka} & \textbf{Body / parametry} & \textbf{Opis} \\
\midrule
\endhead
\httpmethod{POST} & \ep{/admin/file/find} & \ep{FilesFilterCriteria} &
  Wyszukiwanie plików wg filtrów, sortowania i paginacji. \\
\httpmethod{GET} & \ep{/admin/file/download/\{actionId\}} & \ep{actionId} (UUID, ścieżka) &
  Pobranie zawartości pliku powiązanego z danym wpisem logu akcji. \\
\httpmethod{GET} & \ep{/admin/actionLogs/processNames} & -- &
  Lista unikalnych nazw procesów występujących w logach akcji. \\
\httpmethod{POST} & \ep{/admin/fileAlert/find} & \ep{FileAlertsFilterCriteria} &
  Wyszukiwanie alertów plikowych (wraz z metadanymi powiązanego pliku). \\
\httpmethod{GET} & \ep{/admin/fileAlert/\{guid\}} & \ep{guid} (ścieżka) &
  Szczegóły pojedynczego alertu plikowego. \\
\httpmethod{POST} & \ep{/admin/fileAlert/stats} & \ep{FileAlertStatsParams} &
  Statystyki alertów (np. liczba wg poziomu istotności w przedziale czasowym). \\
\httpmethod{POST} & \ep{/admin/user/find} & \ep{UsersFilterCriteria} &
  Wyszukiwanie użytkowników zarejestrowanych w systemie. \\
\httpmethod{POST} & \ep{/admin/actionLogs/find} & \ep{ActionLogFilterCriteria} &
  Wyszukiwanie logów akcji (dostępne wyłącznie w profilu \emph{collector}). \\
\httpmethod{GET} & \ep{/admin/actionLog/\{guid\}} & \ep{guid} (ścieżka) &
  Szczegóły pojedynczego wpisu logu akcji. \\
\httpmethod{POST} & \ep{/admin/group/find} & \ep{GroupFilterCriteria} &
  Wyszukiwanie grup użytkowników (profil \emph{collector}). \\
\httpmethod{POST} & \ep{/admin/group} & \ep{GroupDTO\{name, description, users[]\}} &
  Utworzenie lub aktualizacja grupy użytkowników. \\
\httpmethod{DELETE} & \ep{/admin/group} & \ep{GroupNameDTO\{name\}} &
  Usunięcie grupy. \\
\httpmethod{POST} & \ep{/admin/group/info} & \ep{GroupNameDTO\{name\}} &
  Szczegóły grupy wraz z listą przypisanych użytkowników. \\
\httpmethod{GET} & \ep{/admin/reports} & -- &
  Lista zdefiniowanych raportów. \\
\httpmethod{GET} & \ep{/admin/reports/\{guid\}} & \ep{guid} (ścieżka) &
  Szczegóły raportu, w tym lista jego parametrów wejściowych. \\
\httpmethod{POST} & \ep{/admin/reports/run/\{guid\}} & \ep{guid} (ścieżka), \ep{Map<String,Object>} (parametry) &
  Wykonanie zapisanego w bazie zapytania SQL raportu z podanymi parametrami. \\
\httpmethod{POST} & \ep{/admin/rule/find} & \ep{RuleFilterCriteria} &
  Wyszukiwanie reguł DLP. \\
\httpmethod{DELETE} & \ep{/admin/rule/\{guid\}} & \ep{guid} (ścieżka) &
  Usunięcie reguły DLP. \\
\httpmethod{POST} & \ep{/admin/rule} & \ep{multipart/form-data}: część \ep{file} (plik wtyczki .py, opcjonalny), część \ep{metadata} (\ep{RuleDTO}) &
  Utworzenie lub aktualizacja reguły DLP (nazwa, wyrażenie regularne, definicja, wtyczka, flaga aktywności). \\
\bottomrule
\caption{Endpointy \code{AdminController}.}
\end{longtable}
\endgroup

\subhead{Device API -- \texorpdfstring{\code{/device}}{/device}}

Interfejs pobierania konfiguracji agenta. Wymaga roli \ep{USER}, przyznawanej po pozytywnej
weryfikacji nagłówków \ep{WinKey}/\ep{BiosSN}. Aktywny w profilu \emph{collector}.

\begingroup\footnotesize
\begin{longtable}{M{1.4cm} M{2.8cm} M{5.2cm} M{5.0cm}}
\toprule
\rowcolor{tablehead!15}
\textbf{Metoda} & \textbf{Ścieżka} & \textbf{Nagłówki / Body} & \textbf{Opis} \\
\midrule
\endhead
\httpmethod{POST} & \ep{/device/getConfig} &
  Nagłówki: \ep{WinKey}, \ep{BiosSN}. Body: \ep{UserIdentifier\{userName, ip, winKey, biosSN\}} &
  Zwraca aktualną konfigurację agenta (\ep{DeviceConfigDTO}: interwał synchronizacji, maksymalny
  rozmiar pliku, lista monitorowanych rozszerzeń, lista wykluczonych procesów/ścieżek, flaga
  wysyłania plików). Jeśli użytkownik/stacja o danym \ep{WinKey}/\ep{BiosSN} nie istnieje --
  zostaje utworzony automatycznie. \\
\bottomrule
\caption{Endpoint \code{DeviceController}.}
\end{longtable}
\endgroup

\subhead{Log API -- \texorpdfstring{\code{/log}}{/log}}

Interfejs przyjmowania logów aktywności z agenta. Wymaga roli \ep{USER}. Aktywny w profilu
\emph{collector}.

\begingroup\footnotesize
\begin{longtable}{M{1.4cm} M{2.8cm} M{5.2cm} M{5.0cm}}
\toprule
\rowcolor{tablehead!15}
\textbf{Metoda} & \textbf{Ścieżka} & \textbf{Nagłówki / Body} & \textbf{Opis} \\
\midrule
\endhead
\httpmethod{POST} & \ep{/log/test/} & \ep{String} &
  Endpoint testowy (echo), bez wymaganej autoryzacji -- służy do weryfikacji łączności agenta z serwerem. \\
\httpmethod{POST} & \ep{/log/actionLogs} &
  Nagłówki: \ep{WinKey}, \ep{BiosSN}. Body: \ep{ActionLogUploadDTO\{userName, ip, winKey, biosSN, agentVersion, logs[]\}} &
  Zbiorczy import logów operacji na plikach zebranych przez agenta; każdy wpis otrzymuje status
  początkowy \ep{WAITING} (oczekuje na sprawdzenie przez parser). \\
\httpmethod{POST} & \ep{/log/fileActions} &
  Nagłówki: \ep{WinKey}, \ep{BiosSN}. Body: \ep{FileActionUploadDTO\{..., fileActions[]\}} &
  Zapis powiązań pomiędzy wpisem akcji a przesłanym plikiem. \\
\bottomrule
\caption{Endpointy \code{LogController}.}
\end{longtable}
\endgroup

\subhead{File Upload API -- \texorpdfstring{\code{/upload}}{/upload}}

\begingroup\footnotesize
\begin{longtable}{M{1.4cm} M{2.8cm} M{5.2cm} M{5.0cm}}
\toprule
\rowcolor{tablehead!15}
\textbf{Metoda} & \textbf{Ścieżka} & \textbf{Nagłówki / Body} & \textbf{Opis} \\
\midrule
\endhead
\httpmethod{POST} & \ep{/upload/file} &
  Nagłówki: \ep{WinKey}, \ep{BiosSN}. \ep{multipart/form-data}: część \ep{file} (opcjonalna),
  część \ep{metadata} (\ep{FileUploadDTO}) &
  Przyjmuje plik skopiowany przez użytkownika (np. wynik operacji schowka lub kopiowania pliku)
  do dalszej analizy. Deduplikacja po skrócie SHA-256 pliku (\ep{fileHash}) -- ten sam plik nie
  jest zapisywany dwukrotnie. Plik zapisywany jest fizycznie na dysku serwera. Zwraca \ep{guid}
  utworzonego rekordu pliku. \\
\bottomrule
\caption{Endpoint \code{FileUploadController}.}
\end{longtable}
\endgroup

\subhead{Alert API -- \texorpdfstring{\code{/alert}}{/alert} (profil \emph{aggregator})}

Interfejs przyjmujący zbiorczo alerty wysyłane przez serwery typu \emph{collector} do centralnego
serwera \emph{aggregator}. Wymaga roli \ep{DLP\_COLLECTOR}, przyznawanej na podstawie tokenu
klienta.

\begingroup\footnotesize
\begin{longtable}{M{1.4cm} M{2.4cm} M{5.6cm} M{5.0cm}}
\toprule
\rowcolor{tablehead!15}
\textbf{Metoda} & \textbf{Ścieżka} & \textbf{Body} & \textbf{Opis} \\
\midrule
\endhead
\httpmethod{POST} & \ep{/alert/add} & \ep{List<FileAlertInfoDTO>} &
  Odbiera paczkę alertów od serwera typu \emph{collector}. Dla każdego alertu wyszukiwany lub
  tworzony jest odpowiadający mu plik (po \ep{guid}); alerty o już istniejącym \ep{guid} są
  pomijane (idempotentność). \\
\bottomrule
\caption{Endpoint \code{AlertController}.}
\end{longtable}
\endgroup

\subhead{API wewnętrzne (\texorpdfstring{\code{/internal/...}}{/internal/...})}
\label{sec:internal-api}

Interfejs przeznaczony wyłącznie dla lokalnego procesu \ep{dlp-parser}. Dostępny tylko na
osobnym porcie wewnętrznym serwera (domyślnie \ep{8081}) -- próba wywołania tych ścieżek na
porcie głównym kończy się błędem \ep{RESOURCE\_NOT\_FOUND}. Ze względu na izolację sieciową
endpointy te nie wymagają dodatkowego uwierzytelniania aplikacyjnego (\ep{permitAll}).

\begingroup\footnotesize
\begin{longtable}{M{1.4cm} M{4.6cm} M{4.5cm} M{4.0cm}}
\toprule
\rowcolor{tablehead!15}
\textbf{Metoda} & \textbf{Ścieżka} & \textbf{Body / parametry} & \textbf{Opis} \\
\midrule
\endhead
\httpmethod{GET} & \ep{/internal/actionLog/nextToCheck} & \ep{pageSize} (domyślnie 5) &
  Zwraca kolejną paczkę logów akcji oczekujących na sprawdzenie przez parser. \\
\httpmethod{GET} & \ep{/internal/file/download/\{guid\}} & \ep{guid} (identyfikator akcji, ścieżka) &
  Pobranie surowej zawartości binarnej pliku. \\
\httpmethod{GET} & \ep{/internal/file/nextToCheck/\{ruleGuid\}} & \ep{ruleGuid} (ścieżka) &
  Zwraca kolejny plik do sprawdzenia wskazaną regułą DLP. \\
\httpmethod{POST} & \ep{/internal/file/updateRuleExecutionResult} & \ep{RuleExecDTO\{ruleGuid, fileGuid, status, alerts[]\}} &
  Zapisuje wynik wykonania reguły na danym pliku wraz z ewentualnie wygenerowanymi alertami. \\
\httpmethod{POST} & \ep{/internal/file/update} & \ep{FileProcessingUpdateDTO\{guid, status, alerts[]\}} &
  Aktualizuje status przetwarzania pliku (np. na \ep{COMPLETED}) oraz zapisuje wygenerowane alerty. \\
\httpmethod{GET} & \ep{/internal/file/nextToParse/\{status\}} & \ep{status} (\ep{FileProcessingStatus}, ścieżka) &
  Zwraca surową zawartość kolejnego pliku oczekującego na parsowanie w danym statusie. \\
\httpmethod{GET} & \ep{/internal/file/nextToParse/content/\{status\}} & \ep{status} (ścieżka) &
  Zwraca metadane i wcześniej wyekstrahowaną treść tekstową pliku (bez surowych bajtów). \\
\httpmethod{POST} & \ep{/internal/file/saveParsedContent} & \ep{FileProcessingUpdateDTO\{guid, status, content\}} &
  Zapisuje wyekstrahowaną treść tekstową pliku (np. wynik OCR) oraz zmienia jego status na
  \ep{READY\_TO\_PARSE}. \\
\httpmethod{POST} & \ep{/internal/group/import} & \ep{ImportGroupsDTO\{groups[]\}} &
  Pełna synchronizacja grup i przypisań użytkowników na podstawie danych z Active Directory
  (dodaje/aktualizuje/usuwa grupy zgodnie z przekazanym stanem docelowym). \\
\bottomrule
\caption{Endpointy wewnętrzne (\code{ActionLogInternalController}, \code{FileInternalController}, \code{GroupInternalController}).}
\end{longtable}
\endgroup

\noindent\textbf{Uwaga:} metoda aktualizacji logów akcji (\ep{/internal/actionLog/update}) jest
obecnie zakomentowana w kodzie serwera i nieaktywna, mimo że widnieje w kolekcji testowej Postman
-- traktować jako element planowany, a nie wdrożony.

% =========================================================================
\subsubsection*{3.5. Szczegółowy opis funkcji -- DLP Agent (C\#)}
% =========================================================================

\subhead{Architektura i uruchomienie}

\ep{DLPAgent} jest usługą systemu Windows uruchamianą przez framework \textbf{Topshelf}. Punktem
wejścia jest metoda \ep{Program.Main(string[] args)}:

\begin{itemize}
  \item tryb \ep{--console} -- uruchomienie interaktywne w konsoli, przeznaczone do debugowania;
  \item tryb domyślny -- rejestracja jako usługa Windows (\ep{HostFactory.Run}), z serwisem
        \ep{DLPMonitor} obsługującym zdarzenia \ep{Start}/\ep{Stop}.
\end{itemize}

Przy starcie usługa pobiera numer seryjny BIOS oraz klucz produktu Windows (funkcje
\ep{Tools.GetBiosSerialNumber()}, \ep{Tools.GetWindowsProductKey()}) -- wartości te służą jako
trwały identyfikator stacji roboczej i podstawa uwierzytelniania wobec serwera. Następnie odczytuje
minimalną konfigurację z rejestru systemowego (\ep{ConfManager.GetConfigurationFromRegistry()}),
weryfikuje łączność z serwerem i pobiera z niego aktualną, pełną konfigurację
(\ep{ConfManager.CheckForUpdate()}).

Solucja \ep{DLPAgent.sln} obejmuje trzy projekty: \ep{DLPAgent} (właściwa usługa),
\ep{DLPAgentSetup} (instalator MSI) oraz \ep{DLPClipboardMonitor} (aplikacja WinForms
monitorująca schowek). Równolegle rozwijany jest niezależny projekt
\ep{DLPClipboardMonitor\_net80} -- migracja monitora schowka do platformy .NET 8, niepodpięta
jeszcze do głównej solucji ani instalatora.

\subhead{Monitorowanie operacji plikowych}

Klasa \ep{services/FileServiceMonitor.cs} realizuje monitoring plikowy w oparciu o mechanizm ETW
(\emph{Event Tracing for Windows}):

\begin{itemize}[leftmargin=1.6em]
  \item \ep{StartMonitoring()} -- tworzy sesję ETW \ep{TraceEventSession("FileMonitorSession")},
        włącza dostawcę jądra (\emph{kernel provider}) dla zdarzeń \ep{FileIOInit}, \ep{FileIO}
        oraz \ep{Process}, a także dostawcę \ep{TerminalServices} (śledzenie sesji), po czym
        podpina odpowiednie procedury obsługi zdarzeń.
  \item \ep{StopMonitoring()} -- zatrzymuje sesję ETW.
  \item Procedury obsługi zdarzeń: \ep{OnDynamicEvent} (logowanie/wylogowanie/reconnect sesji),
        \ep{OnProcessStart}, \ep{OnProcessStop}, \ep{OnFileCreate}, \ep{OnFileRead},
        \ep{OnFileWrite}, \ep{OnFileDelete}, \ep{OnFileRename}, \ep{OnFileCleanup}.
  \item \ep{ClassifyAndFilter(pid, fileName)} -- odrzuca zdarzenia dotyczące plików/procesów
        systemowych i tła (ścieżki zawierające katalogi \code{windows} i \code{temp}, procesy typu
        \ep{SearchIndexer}, \ep{MsMpEng}, \ep{svchost}), procesy nieinteraktywne oraz sesje
        headless.
  \item \ep{LogFile(...)} -- dla zdarzeń przechodzących przez filtr: liczy skrót SHA-256 pliku
        (\ep{HashStorage.GetFileHash}), dodaje wpis do wspólnej listy logów akcji, a jeżeli plik
        spełnia kryteria wysyłki (\ep{Filter.FileUploadFilter}), zleca jego skopiowanie
        w tle (\ep{SyncFileHelper.BackgroundCopyWhenIdle}).
\end{itemize}

Klasyfikacja operacji użytkownika na wyższym poziomie (np. rozróżnienie \emph{kopiuj/wklej},
\emph{archiwizacja}) jest zamodelowana w postaci wyliczenia \ep{UserActionType}, lecz -- zgodnie
z komentarzami w kodzie (\ep{HandleUserFileAccess}) -- sama logika korelacji zdarzeń w spójne
akcje użytkownika nie jest jeszcze zaimplementowana.

\subhead{Orkiestracja usługi}

Klasa \ep{services/DLPMonitor.cs} spina wszystkie komponenty w jeden cykl życia usługi:

\begin{itemize}[leftmargin=1.6em]
  \item \ep{Start(hostControl)} -- włącza wymagane uprawnienia systemowe
        (\ep{ServiceHelper.EnablePrivilege}, w tym \ep{SeAssignPrimaryTokenPrivilege},
        \ep{SeTcbPrivilege} potrzebne do uruchamiania procesów w cudzej sesji), uruchamia monitor
        plikowy (\ep{FileServiceMonitor.StartMonitoring}), opcjonalnie monitor sesji użytkownika
        (\ep{UserSessionMonitor.StartMonitoring}) oraz cykliczny \ep{Timer} odpalany co
        \ep{ConfManager.timeInterval} sekund.
  \item \ep{Stop(hostControl)} -- zatrzymuje wszystkie powyższe komponenty.
  \item \ep{onTimer(state)} -- obsługa cyklu synchronizacji: wywołanie
        \ep{Synchronizer.SyncData}, sprawdzenie łączności z serwerem, ewentualne odświeżenie
        konfiguracji (\ep{ConfManager.CheckForUpdate}) oraz -- jeśli włączona jest flaga
        \ep{sendFiles} -- wysyłka zgromadzonych plików (\ep{SyncFileHelper.SendFiles}).
\end{itemize}

\subhead{Monitorowanie sesji i uruchamianie monitora schowka}

Klasa \ep{services/UserSessionMonitor.cs} odpowiada za wykrywanie aktywnych sesji użytkowników
Windows (lokalnych i terminalowych) i uruchamianie w ich kontekście procesu monitorującego schowek:

\begin{itemize}[leftmargin=1.6em]
  \item \ep{StartMonitoring()}/\ep{StopMonitoring()} -- pętla sprawdzająca sesje co 1 sekundę.
  \item \ep{OnUserLogon(sessionId)}/\ep{OnUserLogoff(sessionId)} -- reakcja na zmiany stanu sesji.
  \item \ep{LaunchClipboardMonitor(sessionId)} -- za pomocą wywołań natywnych API Windows
        (\ep{WTSQueryUserToken}, \ep{CreateEnvironmentBlock}, \ep{CreateProcessAsUser})
        uruchamia proces \ep{DLPClipboardMonitor.exe} bezpośrednio na pulpicie zalogowanego
        użytkownika (stacja \ep{winsta0}, pulpit \ep{default}), co jest niezbędne, ponieważ usługa
        Windows działa w izolowanej sesji 0 i nie ma bezpośredniego dostępu do schowka użytkownika.
  \item \ep{CheckExistingSessions()} -- weryfikuje, czy proces monitora schowka nadal działa dla
        każdej aktywnej sesji, i restartuje go w razie potrzeby.
\end{itemize}

\subhead{Monitorowanie schowka (DLPClipboardMonitor)}

Osobna aplikacja \ep{DLPClipboardMonitor.exe} (WinForms, klasa
\ep{ClipboardMonitor : NativeWindow, IDisposable}) rejestruje się jako nasłuchujący zmian schowka
(\ep{AddClipboardFormatListener}) i reaguje na komunikat systemowy \ep{WM\_CLIPBOARDUPDATE}
w nadpisanej metodzie \ep{WndProc}. Zastosowano mechanizm "uspokajający" (timer 350~ms), aby
uniknąć wielokrotnego zrzucania tej samej zawartości przy jednej operacji kopiowania.

Metoda \ep{GetClipboardContent} odczytuje bieżącą zawartość schowka
(\ep{Clipboard.GetDataObject()}) i dla rozpoznanych formatów (tekst, HTML, mapa bitowa, RTF,
lista plików) zapisuje ją do plików w lokalnym katalogu cache. Dla nierozpoznanych strumieni danych
rozszerzenie pliku jest wykrywane na podstawie sygnatury bajtowej (\ep{Tools.DetectExtension}).
W wersji .NET Framework dodatkowo filtrowane są zrzuty formatów bitmapowych pochodzących z aplikacji
pakietu biurowego (Excel, Word, PowerPoint, Outlook), aby ograniczyć ilość generowanego szumu.

\subhead{Komunikacja z serwerem backendowym}

Cała komunikacja HTTP odbywa się za pośrednictwem klasy \ep{utils/Http.cs}
(\ep{SendRequest(url, content)}), która do każdego żądania dołącza nagłówki \ep{WinKey},
\ep{BiosSN} oraz \ep{AuthToken} (skrót SHA-256 z kombinacji klucza Windows, numeru seryjnego
BIOS i bieżącej daty). Agent wywołuje trzy grupy endpointów serwera opisane szczegółowo w
podrozdziale~\ref{sec:dlp-api}: pobranie konfiguracji (\ep{/device/getConfig}), wysyłkę logów akcji
(\ep{/log/actionLogs}) oraz wysyłkę skopiowanych plików wraz z metadanymi
(\ep{/upload/file}, \ep{multipart/form-data}).

\subhead{Konfiguracja i przechowywanie stanu}

Klasa \ep{utils/ConfManager.cs} zarządza konfiguracją agenta w dwóch warstwach:

\begin{itemize}[leftmargin=1.6em]
  \item lokalny rejestr Windows, klucz \code{HKLM/SOFTWARE/Atende/DLPAgent}
        (odczyt metodą \code{GetConfigurationFromRegistry}, zapis metodą \code{SaveConfigurationToRegistry});
  \item konfiguracja centralna pobierana z serwera metodą \ep{CheckForUpdate()} (wywołanie
        \ep{/device/getConfig}), która nadpisuje wartości lokalne przy każdym starcie usługi oraz
        cyklicznie co \ep{TimeInterval} sekund.
\end{itemize}

Domyślne wartości instalacyjne (plik \ep{resources/DLPAgent.reg}) obejmują m.in.: adres serwera
(\ep{Host}), maksymalny rozmiar wysyłanego pliku (\ep{MaxFileSize}, domyślnie 5~MB), interwał
synchronizacji (\ep{TimeInterval}, domyślnie 30~s), listę monitorowanych rozszerzeń plików
(m.in. \ep{doc, docx, xls, xlsx, pdf, ppt, pptx, odt, csv}) oraz listę procesów wykluczonych
z monitorowania.

\subhead{Logowanie}

Zdarzenia diagnostyczne i błędy zapisywane są do \textbf{Dziennika zdarzeń systemu Windows}
(\emph{Windows Event Log}) poprzez metody \ep{Tools.LogInfo}/\ep{Tools.LogError}, pod źródłem
odpowiadającym nazwie usługi (\ep{AtendeDLPMonitor}). Niezależnie od tego, klasa
\ep{services/EtwCsvLogger.cs} umożliwia -- w celach diagnostycznych -- zapis surowych zdarzeń ETW
do pliku CSV w tle, z wykorzystaniem kolejki \ep{BlockingCollection<string>}.

% =========================================================================
\subsubsection*{3.6. Szczegółowy opis funkcji -- DLP Backend Server (Java)}
% =========================================================================

\subhead{Architektura i profile pracy}

\ep{dlp-backend-server} zbudowany jest w oparciu o Spring Boot 3.4 (Java 17), z wykorzystaniem
\textbf{Spring Data JDBC} (nie JPA/Hibernate) w połączeniu z QueryDSL do budowy zapytań
filtrujących, Spring Security do uwierzytelniania, Spring Cloud OpenFeign do komunikacji między
instancjami serwera oraz MapStruct do mapowania encji na DTO. Baza danych: MySQL/MariaDB.

Ta sama baza kodu, w zależności od aktywnego profilu Spring (\ep{--spring.profiles.active=...}),
pełni jedną z dwóch ról:

\begin{itemize}[leftmargin=1.6em]
  \item \textbf{\ep{collector}} -- instancja wdrażana lokalnie u klienta. Uruchamia dodatkowy,
        wewnętrzny serwer HTTP (\ep{DlpInternalHttpServer}) na osobnym porcie, przeznaczony do
        komunikacji z lokalnym procesem \ep{dlp-parser}; przyjmuje dane od agentów; wysyła
        zagregowane alerty do serwera typu \emph{aggregator} przy pomocy klienta Feign
        (\ep{FileAlertsNotifier}, zadanie cykliczne).
  \item \textbf{\ep{aggregator}} -- centralna instancja zbierająca alerty z wielu wdrożeń typu
        \emph{collector}, uwierzytelniająca klientów tokenem przypisanym do rekordu \ep{Customer},
        uruchamiana z włączonym SSL.
\end{itemize}

Kontrolery API panelu administracyjnego (\ep{/auth}, \ep{/admin}) są aktywne niezależnie od profilu.

\subhead{Struktura pakietów}

\begin{itemize}[leftmargin=1.6em]
  \item \ep{auth/} -- konfiguracja Spring Security, generowanie i walidacja JWT
        (\ep{auth/jwt}), filtr tokenu klienta dla profilu \emph{aggregator}
        (\ep{auth/token}), filtr \ep{WinKey} dla agentów (\ep{auth/winkey}).
  \item \ep{client/} -- klient Feign do komunikacji z serwerem \emph{aggregator}, komponent
        wysyłający alerty (\ep{FileAlertsNotifier}), dekoder błędów REST.
  \item \ep{config/} -- właściwości konfiguracyjne aplikacji, konfiguracja harmonogramu zadań.
  \item \ep{controller/} -- osiem kontrolerów REST opisanych w podrozdziale~\ref{sec:dlp-api}.
  \item \ep{db/model} -- encje Spring Data JDBC; \ep{db/repository} -- interfejsy repozytoriów
        (Spring Data JDBC/QueryDSL); \ep{DlpDBConnector}/\ep{DlpJDBCConnector} -- centralna
        klasa łącząca logikę biznesową z dostępem do danych.
  \item \ep{error/} -- hierarchia wyjątków domenowych, \ep{ApiErrorCode}, globalny handler
        błędów (\ep{RestExceptionHandler}).
  \item \ep{filter/} -- filtr izolujący endpointy wewnętrzne na osobnym porcie.
  \item \ep{mail/} -- wysyłka e-maili (reset hasła) przez SMTP (SendGrid).
  \item \ep{model/} i \ep{model/dto} -- klasy odpowiedzi API (\ep{*ApiResult}) oraz DTO
        żądań/odpowiedzi, kryteriów filtrowania i grup walidacji.
  \item \ep{schedule/} -- \ep{ScriptScheduler}, cykliczne uruchamianie skryptów silnika parsera.
  \item \ep{services/} -- \ep{LogWriter} (zapis logów alertów).
  \item \ep{storage/} -- przechowywanie plików i wtyczek reguł na lokalnym dysku.
\end{itemize}

Warto odnotować, że projekt nie stosuje klasycznej, oddzielnej warstwy \emph{Service} -- kontrolery
komunikują się bezpośrednio z klasą \ep{DlpJDBCConnector}, która łączy w sobie logikę biznesową
i dostęp do danych (wzorzec \emph{fat connector/repository}).

\subhead{Model danych}

Najważniejsze encje przechowywane w bazie danych (nazwy tabel w nawiasach):

\begin{itemize}[leftmargin=1.6em]
  \item \ep{File} (\ep{dlp\_file}) -- metadane i (opcjonalnie) wyekstrahowana treść pliku:
        \ep{guid}, \ep{fileHash}, \ep{fileName}, \ep{uploadDate}, \ep{contentType},
        \ep{fileSize}, \ep{content}, \ep{status} (\ep{FileProcessingStatus}).
  \item \ep{FileAlert} (\ep{dlp\_file\_alert}) -- alert wykryty dla pliku lub logu akcji:
        \ep{guid}, \ep{type} (\ep{AlertType}), \ep{severity} (\ep{AlertSeverity}),
        \ep{alertSource} (plik / log akcji), \ep{ruleGuid}, \ep{status} (\ep{AlertStatus}),
        \ep{createdDate}, \ep{sentDate}.
  \item \ep{ActionLog} (\ep{dlp\_action\_log}) -- pojedyncza operacja użytkownika na pliku:
        \ep{guid}, \ep{userId}, \ep{fileName}, \ep{filePath}, \ep{operationType}
        (\ep{OperationType}: \ep{CREATE}, \ep{VIEW}, \ep{MODIFY}, \ep{DELETE},
        \ep{RENAME}, \ep{COPY}, \ep{PASTE}, \ep{CLIPBOARD\_COPY}, \ep{CLIPBOARD\_PASTE}),
        \ep{processId/Name/Owner}, \ep{ip}, \ep{winKey}, \ep{biosSN}, \ep{agentVersion}.
  \item \ep{Rule} (\ep{dlp\_rule}) -- definicja reguły DLP: \ep{guid}, \ep{name},
        \ep{regex}, \ep{definition}, \ep{plugin} (nazwa pliku wtyczki Python),
        \ep{enabled}.
  \item \ep{User} (\ep{dlp\_user}) -- reprezentacja stacji roboczej/użytkownika: \ep{guid},
        \ep{customerId}, \ep{winKey}, \ep{biosSN}, \ep{lastSeen}.
  \item \ep{Group}/\ep{UserGroup} (\ep{dlp\_group}) -- grupy użytkowników (synchronizowane
        z Active Directory) i ich przypisania.
  \item \ep{Admin} (\ep{dlp\_admin}) -- konto administratora panelu (login, e-mail,
        zahaszowane hasło).
  \item \ep{AdminSession}, \ep{RevokedAccessToken}, \ep{ConfirmationToken} -- zarządzanie
        sesjami JWT i tokenami resetu hasła.
  \item \ep{Customer} (\ep{dlp\_customer}) -- rekord reprezentujący instalację typu
        \emph{collector} w kontekście serwera \emph{aggregator}, zawiera token autoryzacyjny.
  \item \ep{Report}/\ep{ReportParameter} (\ep{dlp\_report}) -- zdefiniowane raporty
        (zapytanie SQL wraz z parametrami wejściowymi).
\end{itemize}

\subhead{Kluczowe klasy serwisowe}

\begin{itemize}[leftmargin=1.6em]
  \item \textbf{\ep{DlpDBConnector} / \ep{DlpJDBCConnector}} -- centralny komponent łączący
        logikę biznesową z dostępem do danych. Grupy metod: obsługa plików
        (\ep{findFiles}, \ep{saveFile}, \ep{findOrCreateFile}, \ep{getNextFileToCheck/Parse}),
        obsługa alertów (\ep{saveFileAlerts}, \ep{findFileAlerts}, \ep{getFileAlertStats},
        \ep{getNotSentFileAlerts}, \ep{updateFileAlertsStatus}), obsługa logów akcji
        (\ep{saveActionLogs}, \ep{findActionLogs}, \ep{findActionLogsToCheck}), obsługa
        użytkowników i grup (\ep{getUser} -- tworzy użytkownika, jeśli nie istnieje,
        \ep{updateGroups} -- pełna synchronizacja z AD), obsługa reguł
        (\ep{findRules}, \ep{saveRule}, \ep{deleteRule}, \ep{saveRuleExecutionResult}),
        obsługa sesji administratora (\ep{saveAdminSession}, \ep{logoutAllAdminSessions},
        \ep{isAccessTokenRevoked}), obsługa raportów (\ep{findAllReports}, \ep{executeReport}
        -- dynamiczne wykonanie zapisanego zapytania SQL z podstawionymi parametrami).
  \item \textbf{\ep{JwtUtil}} -- generowanie, parsowanie i walidacja tokenów JWT (algorytm
        HS512), tworzenie pary tokenów dostęp/odświeżanie wraz z zapisem sesji administratora.
  \item \textbf{\ep{FileAlertsNotifier}} (profil \emph{collector}) -- zadanie cykliczne
        (\ep{@Scheduled}) wysyłające do serwera \emph{aggregator} paczki (po 50) alertów jeszcze
        niewysłanych, aktualizujące status wysyłki (\ep{SUCCESSFULLY\_SENT}/\ep{SENDING\_FAILED}).
  \item \textbf{\ep{ScriptScheduler}} -- cykliczne (\ep{@Scheduled(fixedRateString = "\$\{python.fixedRate\}")})
        uruchamianie skryptów Pythona silnika \ep{dlp-parser} jako procesów zewnętrznych,
        zabezpieczone blokadą (\ep{ReentrantLock}) uniemożliwiającą równoległe uruchomienie
        tego samego skryptu.
  \item \textbf{\ep{LocalStorageConnector} / \ep{FileStorageConnector}} -- przechowywanie
        przesłanych plików (\ep{\{guid\}.file}) oraz wtyczek reguł (\ep{\{guid\}.py}) na
        lokalnym dysku serwera.
  \item \textbf{\ep{DlpDbDTOMapper}} (MapStruct) -- mapowanie pomiędzy encjami bazodanowymi
        a obiektami DTO wymienianymi przez API.
  \item \textbf{\ep{RestExceptionHandler}} -- globalny handler wyjątków (\ep{@ControllerAdvice}),
        tłumaczący wyjątki domenowe i techniczne na spójne odpowiedzi \ep{ApiResult} wraz
        z odpowiednim kodem błędu i statusem HTTP.
  \item \textbf{\ep{MailSender}/\ep{SMTPMailSender}} -- wysyłka e-maili transakcyjnych
        (reset hasła) przez SendGrid SMTP.
\end{itemize}

\subhead{Harmonogram zadań}

Serwer backendowy pełni rolę \emph{orkiestratora} dla silnika parsującego: klasa
\ep{ScriptScheduler} cyklicznie uruchamia skrypty Pythona z katalogu \ep{dlp-parser/scripts}
jako niezależne procesy systemowe, z ustalonym limitem czasu wykonania. Wymiana danych pomiędzy
serwerem a uruchomionym skryptem odbywa się w całości poprzez API wewnętrzne opisane w
podrozdziale~\ref{sec:internal-api} -- skrypt, mimo że uruchamiany lokalnie, komunikuje się
z serwerem tak samo jak zdalny klient HTTP.

% =========================================================================
\subsubsection*{3.7. Szczegółowy opis funkcji -- DLP Parser (Python)}
% =========================================================================

\subhead{Charakter modułu}

\ep{dlp-parser} nie jest serwerem HTTP -- jest zbiorem niezależnych skryptów wsadowych,
uruchamianych cyklicznie przez \ep{ScriptScheduler} serwera backendowego (rozdz.~4.5). Moduł
pełni funkcję \textbf{klienta} REST API wewnętrznego serwera (\ep{/internal/...}), a jego zadaniem
jest ekstrakcja tekstu z przesłanych plików oraz wykrywanie w nim danych wrażliwych.

\subhead{Punkty wejścia (skrypty)}

\begin{itemize}[leftmargin=1.6em]
  \item \textbf{\ep{scripts/parse\_files.py}} -- pobiera z serwera plik gotowy do parsowania
        (status \ep{READY\_TO\_PARSE}), uruchamia na nim wszystkie zarejestrowane definicje
        parserów o zakresie \ep{sourceTable == "file"}, a wynik (listę alertów) odsyła z powrotem
        do serwera (\ep{mark\_file\_as\_parsed}). Jeśli wśród wygenerowanych alertów znajdzie się
        alert o poziomie \ep{CRITICAL}, wysyłane jest powiadomienie e-mail. Kluczowe funkcje:
        \ep{extract\_data()}, \ep{parse\_files()}, \ep{extract\_critical\_alerts(alerts)}.
  \item \textbf{\ep{scripts/ocr\_file.py}} -- pobiera plik oczekujący na wstępne przetworzenie
        (status \ep{WAITING}), ekstrahuje z niego tekst (w tym poprzez OCR) i odsyła
        wyekstrahowaną treść do serwera (\ep{set\_file\_content}), zmieniając jego status na
        \ep{READY\_TO\_PARSE}. Kluczowa funkcja: \ep{extract\_OCR\_content()}.
  \item \textbf{\ep{scripts/indexer.py}} -- niezależny od pozostałych skrypt indeksujący pliki
        bezpośrednio na dysku stacji, na której jest uruchomiony: skanuje wskazane katalogi, liczy
        skrót SHA-256 każdego pliku, a nowe/zmienione pliki zgłasza do zewnętrznego API
        (\ep{/log/actionLogs}, \ep{/log/fileActions}, \ep{/upload/file}). Stan już zaindeksowanych
        plików przechowuje w lokalnej bazie SQLite.
  \item \textbf{\ep{automate/parser\_generator.py}} -- narzędzie generujące pliki definicji
        parserów JSON na podstawie zawartości monitorowanych katalogów referencyjnych.
\end{itemize}

\subhead{Silnik reguł -- \code{services/parser\_core.py}}

Klasa \ep{ParserCore} realizuje właściwą logikę dopasowywania reguł do treści:

\begin{itemize}[leftmargin=1.6em]
  \item \ep{get\_all\_parsers()} -- listuje wszystkie pliki definicji parserów z katalogu
        wskazanego zmienną środowiskową \ep{PARSERS\_PATH}.
  \item \ep{load\_parser(file\_name)} -- wczytuje plik JSON i waliduje go względem modelu
        \ep{Parser} (Pydantic).
  \item \ep{process\_all(object\_to\_parse)} -- dla każdego zadania (\ep{Task}) zdefiniowanego
        w parserze: buduje tekst wejściowy ze wskazanych pól obiektu (\ep{fieldsToParse}),
        sekwencyjnie stosuje filtry wstępne (\ep{preTasks}, np. lematyzację), a następnie:
        jeśli zadanie zawiera wyrażenie regularne (\ep{regex}) -- dopasowuje je i, przy braku
        wtyczki, samo dopasowanie generuje alert; jeśli zadanie wskazuje wtyczkę (\ep{plugin})
        -- przekazuje do niej dopasowany fragment lub cały tekst w celu dodatkowej walidacji
        i wygenerowania alertu (\ep{severity}, \ep{type}, \ep{taskName}, \ep{parserName}).
\end{itemize}

Definicja parsera (plik JSON) opisuje: nazwę, zakres (\ep{sourceTable}: \ep{file} lub
\ep{actionLog}), rolę (\ep{scope}) oraz listę zadań, z których każde może łączyć wyrażenie
regularne, listę filtrów wstępnych i opcjonalną wtyczkę.

\subhead{Filtry (\code{filters/})}

Filtry implementują wspólny interfejs \ep{Filter} (\ep{set\_text}, \ep{get\_filtered}) i są
stosowane na tekście przed właściwym dopasowaniem reguły:

\begin{itemize}[leftmargin=1.6em]
  \item \ep{LowerCase} -- sprowadzenie tekstu do małych liter.
  \item \ep{RemoveDiactrics} -- usunięcie polskich znaków diakrytycznych.
  \item \ep{Lemmatize} -- lematyzacja tekstu z użyciem modelu spaCy dla języka polskiego.
\end{itemize}

\subhead{Wtyczki (\code{plugins/})}

Wtyczki implementują wspólny interfejs \ep{Plugin} (\ep{set\_text},
\ep{get\_alerts(task, parser\_name)}) i realizują właściwą logikę walidacji danych wrażliwych:

\begin{itemize}[leftmargin=1.6em]
  \item \ep{PeselFind} -- waliduje kandydata jako polski numer PESEL (weryfikacja sumy
        kontrolnej, klasa pomocnicza \ep{services/pesel.py}).
  \item \ep{PhoneNumberFind} -- wykrywa polskie numery telefonów komórkowych (wzorzec
        \ep{+48~XXX~XXX~XXX}).
  \item \ep{AlertFind} -- porównuje nazwę analizowanego pliku z listą plików referencyjnych
        zebranych z monitorowanych katalogów (\ep{automate.ParserGenerator}), generując alert
        przy dopasowaniu podciągu nazwy.
\end{itemize}

Obecnie w repozytorium dostarczona jest jedna gotowa definicja parsera
(\ep{parsers/pesel\_parser.json}, wykrywanie numeru PESEL w treści pliku); wtyczki
\ep{PhoneNumberFind} i \ep{AlertFind} są zaimplementowane, ale wymagają dodania odpowiadającej
im definicji JSON, aby zostały użyte w przebiegu parsowania.

\subhead{Ekstrakcja treści (\code{services/extractor.py})}

Klasa \ep{Extractor} odpowiada za konwersję plików binarnych do postaci tekstowej, obsługując
formaty: \ep{.doc}/\ep{.docx} (konwersja \ep{.doc}$\rightarrow$\ep{.docx} przez LibreOffice,
następnie \ep{python-docx}), \ep{.xls}/\ep{.xlsx} (\ep{xlrd}/\ep{openpyxl}), \ep{.pdf}
(\ep{pypdf}), \ep{.png}/\ep{.jpg} (OCR, \ep{pytesseract}, model językowy polski). Dla
dokumentów Word i PDF opcjonalnie wykonywany jest dodatkowo OCR obrazów osadzonych w treści
dokumentu. Metoda \ep{save\_temporary\_file\_and\_extract\_data} zarządza cyklem życia pliku
tymczasowego (zapis $\rightarrow$ ekstrakcja $\rightarrow$ usunięcie).

\subhead{Komunikacja z serwerem (\code{services/api\_service.py})}

Klasa \ep{ApiService} jest klientem REST API wewnętrznego serwera backendowego (nagłówek
\ep{Authorization: Bearer <token>}, jeśli token jest skonfigurowany). Udostępnia metody
odpowiadające bezpośrednio endpointom opisanym w podrozdziale~\ref{sec:internal-api}:
\ep{get\_file\_metadata\_to\_check()}, \ep{get\_file\_metadata\_to\_parse(status)},
\ep{get\_file\_content\_to\_parse(status)}, \ep{get\_file\_to\_check(guid)},
\ep{mark\_file\_as\_parsed(guid, alerts)}, \ep{mark\_action\_logs\_as\_parsed(action\_log\_alerts)}
oraz \ep{set\_file\_content(guid, file\_content)}.

\subhead{Powiadomienia e-mail (\code{mail/})}

Funkcja \ep{mailer.send\_email(subject, message\_body, logger, critical\_alerts)} wysyła
powiadomienie SMTP (STARTTLS) zawierające opis każdego wykrytego alertu o poziomie \ep{CRITICAL}.
Wywoływana jest z poziomu \ep{scripts/parse\_files.py} po zakończeniu analizy pliku, o ile wśród
wyników znalazł się przynajmniej jeden alert krytyczny.

\subhead{Dane referencyjne i stan lokalny}

Katalog \ep{db/} zawiera statyczne pliki referencyjne (listy popularnych polskich nazwisk wg
płci), obecnie nieużywane przez żaden aktywny filtr lub wtyczkę -- przygotowane pod przyszłą
funkcję wykrywania nazwisk w treści dokumentów. Jedyną rzeczywistą bazą danych utrzymywaną przez
ten moduł jest lokalna baza SQLite tworzona przez \ep{scripts/indexer.py} (tabele \ep{hashes}
-- mapa zaindeksowanych plików po skrócie SHA-256, oraz \ep{meta} -- metadane przebiegu, np.
znacznik czasu ostatniego indeksowania).

% =========================================================================
\subsubsection*{3.8. Bezpieczeństwo i uwierzytelnianie}
\label{ch:bezpieczenstwo}
% =========================================================================

System stosuje \textbf{trzy niezależne mechanizmy uwierzytelniania}, odpowiadające trzem różnym
typom klientów API, egzekwowane przez łańcuch filtrów Spring Security (\ep{OncePerRequestFilter})
oraz adnotacje \ep{@PreAuthorize} na poszczególnych metodach kontrolerów.

\begingroup\footnotesize
\begin{longtable}{M{3.2cm} M{3.6cm} M{3.6cm} M{4.6cm}}
\toprule
\rowcolor{tablehead!15}
\textbf{Strefa API} & \textbf{Klient} & \textbf{Mechanizm} & \textbf{Rola} \\
\midrule
\endhead
\ep{/auth}, \ep{/admin} & Panel administracyjny & JWT (\ep{Authorization: Bearer}), algorytm HS512, z możliwością odwołania tokenu & \ep{ADMIN} \\
\ep{/device}, \ep{/log}, \ep{/upload} & Agent Windows & Nagłówki \ep{WinKey}, \ep{BiosSN}, \ep{AuthToken} (skrót powiązanych wartości) & \ep{USER} \\
\ep{/alert} & Serwer \emph{collector} $\rightarrow$ \emph{aggregator} & Nagłówek \ep{AuthToken} porównywany ze statycznym tokenem klienta (\ep{Customer.authToken}) & \ep{DLP\_COLLECTOR} \\
\ep{/internal/*} & Lokalny proces \ep{dlp-parser} & Izolacja sieciowa -- osobny port serwera, brak uwierzytelniania aplikacyjnego & -- \\
\bottomrule
\caption{Mechanizmy uwierzytelniania w poszczególnych strefach API.}
\end{longtable}
\endgroup

\noindent Dodatkowe elementy modelu bezpieczeństwa:

\begin{itemize}[leftmargin=1.6em]
  \item Hasła administratorów przechowywane są w postaci skrótu \textbf{BCrypt} (współczynnik
        kosztu 13) wraz z dodatkową, tajną wartością (\emph{pepper}); siła nowo ustawianego hasła
        weryfikowana jest biblioteką \textbf{Passay}.
  \item Reset hasła odbywa się poprzez jednorazowy, ograniczony czasowo token
        (\ep{ConfirmationToken}); endpoint \ep{/auth/recoverPassword} zawsze zwraca odpowiedź
        sukcesu, aby nie ujawniać, czy dane konto istnieje w systemie.
  \item API jest bezstanowe (\ep{SessionCreationPolicy.STATELESS}), CSRF jest wyłączone
        (nie dotyczy API bezstanowego), włączone są natomiast nagłówki \ep{X-XSS-Protection}
        oraz polityka \ep{Content-Security-Policy} ograniczająca źródła skryptów.
  \item Konfiguracja CORS jest parametryzowana (\ep{cors.allowedOriginPatterns}) i powinna być
        zawężona do rzeczywistych domen panelu administracyjnego w środowisku produkcyjnym.
  \item Po stronie agenta walidacja certyfikatu SSL serwera może zostać globalnie wyłączona
        (flaga \ep{Globals.certValidation}) -- ustawienie to ma bezpośredni wpływ na
        bezpieczeństwo komunikacji HTTPS i powinno być używane wyłącznie w środowiskach
        testowych.
  \item Uwierzytelnianie agenta (\ep{WinKey}/\ep{BiosSN}/\ep{AuthToken}) nie jest
        standardowym mechanizmem OAuth2/JWT -- opiera się na identyfikatorach sprzętowych stacji
        i wspólnym sekrecie wyliczanym po obu stronach.
\end{itemize}

% =========================================================================
\subsubsection*{3.9. Konfiguracja modułów}
% =========================================================================

\subhead{DLP Agent}

Konfiguracja agenta przechowywana jest w rejestrze systemowym
(\ep{HKLM/SOFTWARE/Atende/DLPAgent}) i nadpisywana
konfiguracją centralną pobieraną z serwera. Najważniejsze parametry:

\begingroup\footnotesize
\begin{longtable}{M{4cm} M{3cm} M{8cm}}
\toprule
\rowcolor{tablehead!15}
\textbf{Parametr} & \textbf{Domyślnie} & \textbf{Znaczenie} \\
\midrule
\endhead
\ep{Host} & \ep{https://<adres>/dlp-srv/} & Bazowy adres URL serwera backendowego. \\
\ep{Enabled} & \ep{1} & Flaga globalnego włączenia monitoringu. \\
\ep{TimeInterval} & 30 s & Okres cyklu synchronizacji logów i odświeżania konfiguracji. \\
\ep{MaxFileSize} & 5 MB & Maksymalny rozmiar pliku kwalifikującego się do wysyłki. \\
\ep{SendFiles} & \ep{0} & Flaga włączenia wysyłki skopiowanych plików (nie tylko logów). \\
\ep{ResendKnownFileInterval} & \ep{0} & Odstęp, po którym już wysłany plik może zostać wysłany ponownie. \\
\ep{Extensions} & lista rozszerzeń dokumentów biurowych & Rozszerzenia plików objęte monitoringiem. \\
\ep{Processes} & lista procesów systemowych i narzędziowych & Procesy wykluczone/uwzględnione w monitoringu (tryb wg \ep{ExcludeIncludeProcessesMode}). \\
\bottomrule
\caption{Kluczowe parametry konfiguracyjne DLP Agenta.}
\end{longtable}
\endgroup

\subhead{DLP Backend Server}

Konfiguracja podzielona jest na plik wspólny (\ep{application.yml}) oraz pliki właściwe dla
profilu (\ep{application-collector.yml}, \ep{application-aggregator.yml}), wybieranego
parametrem \ep{spring.profiles.active}.

\begin{itemize}[leftmargin=1.6em]
  \item \textbf{Wspólne}: klucz szyfrowania konfiguracji (\ep{encrypt.key}), ustawienia CORS,
        konfiguracja szablonu i nadawcy e-maila resetu hasła, konfiguracja SMTP (SendGrid),
        parametry uruchamiania silnika parsującego (\ep{python.*}: interpreter, katalog
        skryptów, limit czasu, częstotliwość harmonogramu).
  \item \textbf{Profil \ep{collector}}: port główny \ep{8080} i port wewnętrzny \ep{8081}
        (z prefiksem \ep{/internal/}), ścieżka pliku domyślnej konfiguracji urządzenia
        (\ep{config.json}), magazyn zaufania SSL do komunikacji z serwerem \emph{aggregator},
        token klienta (\ep{dlp.customer.token}), ustawienia powiadamiania \emph{aggregator}a
        (\ep{dlp.dataAggregator.*}), limit rozmiaru przesyłanych plików (32~MB), parametry
        połączenia z bazą MySQL, ścieżki przechowywania plików i wtyczek reguł, parametry JWT
        (sekret, czas ważności tokenu dostępowego i odświeżającego).
  \item \textbf{Profil \ep{aggregator}}: port \ep{8085} z włączonym SSL (magazyn kluczy
        PKCS\#12), niższy limit rozmiaru przesyłanych plików (2~MB), osobna baza danych, własny
        sekret JWT i nazwa usługi.
\end{itemize}

\subhead{DLP Parser}

Konfiguracja realizowana jest przez plik \ep{.env} (zmienne środowiskowe), odczytywany przy
starcie każdego skryptu:

\begin{itemize}[leftmargin=1.6em]
  \item \ep{SERVER\_HOST}, \ep{INTERNAL\_HOST} -- adresy API serwera backendowego (odpowiednio
        publicznego i wewnętrznego).
  \item \ep{PARSERS\_PATH} -- katalog z plikami definicji parserów JSON.
  \item \ep{STORAGE\_DIR}, \ep{TMP\_DIR}, \ep{RESULTS\_DIR}, \ep{OCR} -- konfiguracja
        katalogów roboczych i włączenia OCR.
  \item \ep{AUTOMATE\_CHECKED\_FOLDERS}, \ep{PARSERS\_OUTPUT\_PATH}, \ep{PARSERS\_TEMPLATE\_PATH}
        -- konfiguracja generatora definicji parserów.
  \item \ep{INDEXER\_EXTENSIONS}, \ep{INDEXER\_SERVER\_URL}, \ep{INDEXER\_DB\_FILE} --
        konfiguracja skryptu indeksującego.
  \item \ep{SPRING\_MAIL\_HOST/PORT/USERNAME/PASSWORD}, \ep{SENDER\_EMAIL}, \ep{RECEIVER\_EMAIL}
        -- konfiguracja SMTP powiadomień o alertach krytycznych.
  \item \ep{LOG\_FILE\_PATH} -- ścieżka pliku logu (rotacja co 5~MB, 3 kopie zapasowe).
\end{itemize}

% % =========================================================================
% \subsubsection*{3.10. Znane ograniczenia i dług techniczny}
% % =========================================================================

% Na potrzeby rzetelności dokumentacji odnotowano poniższe elementy będące w toku rozwoju lub
% wymagające uwagi zespołu utrzymaniowego:

% \begin{itemize}[leftmargin=1.6em]
%   \item \textbf{DLP Agent}: korelacja pojedynczych zdarzeń plikowych w spójne, wysokopoziomowe
%         akcje użytkownika (kopiuj/wklej/archiwizacja) nie jest jeszcze zaimplementowana
%         (\ep{FileServiceMonitor.HandleUserFileAccess}); mechanizm anonimizacji ścieżek plików
%         (\ep{Anonymizer}) istnieje w kodzie, lecz nie jest aktywnie wykorzystywany -- do serwera
%         trafia pełna, nieanonimizowana ścieżka pliku; wersja monitora schowka dla .NET~8
%         (\ep{DLPClipboardMonitor\_net80}) nie jest jeszcze zintegrowana z głównym instalatorem.
%   \item \textbf{DLP Backend Server}: metoda aktualizacji logów akcji przez API wewnętrzne
%         (\ep{/internal/actionLog/update}) jest obecnie zakomentowana w kodzie; kolekcja testowa
%         Postman nie obejmuje wszystkich endpointów zaimplementowanych w kodzie (m.in. zarządzania
%         regułami) i nie powinna być traktowana jako w pełni aktualne źródło prawdy o API.
%   \item \textbf{DLP Parser}: silnik reguł sygnalizowany jest jako niestabilny przy jednoczesnym
%         wczytaniu więcej niż jednej definicji parsera (komentarz w \ep{parser\_core.py});
%         rejestr filtrów (\ep{FILTER\_REGISTRY}) nie obejmuje wszystkich zaimplementowanych klas
%         filtrów (\ep{LowerCase}, \ep{RemoveDiactrics} nie są obecnie dostępne z poziomu
%         definicji JSON); katalog \ep{db/} z listami nazwisk nie jest jeszcze podpięty pod żaden
%         filtr/wtyczkę; plik \ep{README.md} odwołuje się do nieistniejącego skryptu
%         \ep{parse\_action\_logs.py} i wymaga aktualizacji; \ep{Dockerfile} zawiera
%         nieprawidłowe polecenie uruchomieniowe (odwołanie do nieistniejącego modułu
%         \ep{scripts.generate}) i efektywnie służy jedynie jako środowisko developerskie.
% \end{itemize}
